% Self-contained journal-style fragments. Labels and citation keys use q:.
% Required packages: amsmath, amssymb, amsthm, mathtools, hyperref.
% Required theorem environments: theorem, lemma, proposition, corollary, remark.

\section{Completely positive transport and tensor-network gaps}
\label{q:sec-networks}

A transfer operator is usually used to estimate the decay of correlations
along a one-dimensional tensor network.  Here it has a second role: it
controls the angle between two local ground spaces meeting at a vertex
of an arbitrary graph.  The useful comparison is an inequality between
completely positive maps.  Such an inequality remains valid when a wire
is attached to an arbitrarily entangled positive core.  Moreover, two
different edge constraints can interact through that core only when both
of their wire maps depart from their rank-one limits.  This gives a
quadratic, rather than linear, error bound.

The use of transfer operators and ground-space angles for decorated
AKLT models originates in \cite{q:ARLLNY2020}; the extension to bounded-degree
graphs with a sufficient decoration linear in the degree is
\cite{LuciaYoung2023}.  The argument below gives a degree-independent
two-port angle estimate, leading to logarithmic sufficient decoration.
The final conversion from a blocked projector parent to a specified
microscopic Hamiltonian is stated separately, with its normalization
and finite-patch constants explicit.

\subsection{Finite local data}

Let $\mathcal G=(V,E)$ be a finite simple graph without isolated
vertices.  Open regions may also have dangling boundary ports; these
are included in every product over ports below, but carry no internal
edge constraint.  Every incidence $(v,e)$ carries a finite-dimensional exposed
virtual space.  Every edge carries a fixed nonzero vector of full Schmidt
rank on its two exposed spaces.  At each vertex choose an arbitrary
nonzero core map $A_v$, and attach one matrix-product half-wire to
each of its input legs.  Write $T_v$ for the resulting map from the
exposed virtual spaces to the physical block.  The block consists of
the core and its incident half-wires.

Put $C_v=A_v^*A_v$.  Contraction of the half-wires gives the exact identity
\begin{equation}
 T_v^*T_v=\left(\bigotimes_{e\ni v}\Phi_{v,e}\right)(C_v),
 \label{q:gram-channel}
\end{equation}
where each $\Phi_{v,e}$ is a completely positive (CP) map.  One may fix
the convention directly: if $W_{p,a,b}$ is the amplitude of a half-wire
with physical index $p$, core index $a$, and exposed index $b$, its
Kraus matrices are $(K_p)_{b,a}=\overline{W_{p,a,b}}$.  Thus
\[
 (T^*T)_{b',b}
 =\sum_{p,a',a}\overline{W_{p,a',b'}}(C_v)_{a',a}W_{p,a,b}
 =\left(\sum_p K_pC_vK_p^*\right)_{b',b}.
\]
In particular, no partial transpose of an entangled core is involved.

Assume that, for every incidence, there is a rank-one CP map
\[
 \Psi_{v,e}(X)=\operatorname{tr}(\tau_{v,e}X)R_{v,e},
 \qquad \tau_{v,e}>0,\quad R_{v,e}>0,
\]
and a number $0\leq\delta_{v,e}<1$ such that
\begin{equation}
 (1-\delta_{v,e})\Psi_{v,e}
 \ \leq_{\mathrm{CP}}\ \Phi_{v,e}
 \ \leq_{\mathrm{CP}}\ (1+\delta_{v,e})\Psi_{v,e}.
 \label{q:cp-sandwich}
\end{equation}
Here $\Phi\leq_{\mathrm{CP}}\widetilde\Phi$ means that
$\widetilde\Phi-\Phi$ is completely positive.  This is equivalent to a
positive-semidefinite comparison of their Choi matrices
\cite{q:Choi1975}.  The reference
map need not preserve the trace.  A sandwich
$a\Psi\leq_{\mathrm{CP}}\Phi\leq_{\mathrm{CP}}b\Psi$, with $0<a\leq b$,
is centered by replacing $\Psi$ with $(a+b)\Psi/2$ and setting
$\delta=(b-a)/(b+a)$.

Tensor products preserve CP order.  Consequently
\begin{equation}
 T_v^*T_v\geq
 \left[\prod_{e\ni v}(1-\delta_{v,e})\right]
 \operatorname{tr}\!\left[
   \left(\bigotimes_{e\ni v}\tau_{v,e}\right)C_v\right]
 \bigotimes_{e\ni v}R_{v,e}>0.
 \label{q:injectivity}
\end{equation}
Thus every block map is injective, even when the original core has rank
one.  Their tensor product is injective as well: it sends the nonzero
product of internal edge vectors and any nonzero boundary vector to a
nonzero physical state.  This proves frustration-freeness without an
assumption that a general tensor contraction avoids cancellation.
All estimates below are independent of the condition number of $C_v$.

For $e=uv$, contract the fixed edge vector between $T_u$ and $T_v$,
leaving their other exposed legs free, and denote the resulting regional
image by $\mathcal R_e$.  Let $G_e$ be its orthogonal projector in the
two physical blocks, extended by the identity elsewhere.  Define the
norm-one-projector parent
\begin{equation}
 H_{\mathrm{block}}=\sum_{e\in E}(I-G_e).
 \label{q:block-parent}
\end{equation}
Its ground space retains all outer virtual degrees of freedom when the
graph describes an open region.

\begin{theorem}[Two-port CP gap criterion]
\label{q:thm-wire-gap}
Under \eqref{q:cp-sandwich}, set
\begin{equation}
 \mathcal B(x,y)=\frac{2xy(1+xy)}{(1-x)^2(1-y)^2}.
 \label{q:Bxy}
\end{equation}
Let $C$ be the symmetric matrix indexed by edges, with zero diagonal,
with $C_{ef}=0$ for disjoint edges, and with
$C_{ef}=\mathcal B(\delta_{v,e},\delta_{v,f})$ when $e,f$ meet at $v$.
If its spectral radius satisfies $\rho(C)<1$, then
\begin{equation}
 H_{\mathrm{block}}\geq
 [1-\rho(C)](I-G_{\mathrm{all}}),
 \label{q:block-gap}
\end{equation}
where $G_{\mathrm{all}}$ projects onto its full ground space.
The same statement holds uniformly for a family of finite graphs if
their matrices $C$ have a common spectral-radius upper bound below one.
\end{theorem}

The simpler sufficient condition is the row bound

\fig{fig09_quantum_networks}{Completely positive wire estimates give
finite gap certificates. (a) A schematic network has arbitrary nonzero
cores and mixing half-wires; four sites per original edge are drawn only
to illustrate the construction. (b) The exact two-port estimate
$\mathcal B(\delta_i,\delta_j)$ is quadratic at the origin, reflecting
cancellation when either wire is exactly rank one. (c) Substitution of
$\delta_\ell=2/(3^\ell+(-1)^\ell)$ into the degree bound gives
$1-2(\Delta-1)\mathcal B(\delta_\ell,\delta_\ell)$ for a norm-one
blocked parent. Negative lower bounds are plotted as zero, meaning that
this criterion provides no positive certificate there. A microscopic
AKLT gap additionally includes the finite-patch gap and coverage factors
proved in the text.}{fig:networks}
\begin{equation}
 q=\max_{e=uv}\left[
 \sum_{\substack{f\ni u\\f\ne e}}
       \mathcal B(\delta_{u,e},\delta_{u,f})+
 \sum_{\substack{f\ni v\\f\ne e}}
       \mathcal B(\delta_{v,e},\delta_{v,f})\right]<1,
 \label{q:row-criterion}
\end{equation}
which gives gap at least $1-q$.  Rational positive-vector certificates
$Cw\leq rw$, $w>0$, $r<1$, can improve the row bound while avoiding any
floating-point eigenvalue calculation.  Finite edge types suffice when
their incidence counts give a uniform finite comparison matrix.

\subsection{The quadratic cancellation}

For subspaces $M,N$, their Friedrichs cosine is the largest overlap of
unit vectors in $M\ominus(M\cap N)$ and $N\ominus(M\cap N)$.
Removing the entire intersection is essential for degenerate local
ground spaces.

\begin{lemma}[A quotient-metric estimate]
\label{q:lem-schur}
Let $p,q$ be commuting orthogonal projections.  Suppose
\[
 G=I+A+B+D>0,\qquad [A,q]=[B,p]=0,
\]
\[
 \|A\|\leq x,\quad \|B\|\leq y,\quad \|D\|\leq xy,
 \qquad G\geq mI,\quad m=(1-x)(1-y)>0.
\]
The Friedrichs cosine between $G^{1/2}\operatorname{ran}p$ and
$G^{1/2}\operatorname{ran}q$ is at most $\mathcal B(x,y)$.
\end{lemma}

\begin{proof}
Put $W=\operatorname{ran}(pq)$,
$U=\operatorname{ran}(p(I-q))$, and
$V=\operatorname{ran}((I-p)q)$.  They are mutually orthogonal.  The
quotient by the common image of $W$ is represented by minimizing over
the $W$ coordinate in the $G$-metric.  Its $U,V$ cross block is therefore
\[
 S_{UV}=G_{UV}-G_{UW}G_{WW}^{-1}G_{WV}.
\]
Only $W$ is eliminated: the sector
$\operatorname{ran}((I-p)(I-q))$ is not part of the intersection.
The commutation assumptions imply
\[
 G_{UV}=D_{UV},\qquad
 G_{UW}=B_{UW}+D_{UW},\qquad
 G_{WV}=A_{WV}+D_{WV}.
\]
Hence
\[
 \|S_{UV}\|\leq xy+\frac{(y+xy)(x+xy)}m.
\]
For $u\in U$, the diagonal quotient metric obeys
\[
 \inf_{w\in W}\langle u+w,G(u+w)\rangle
 \geq m\inf_{w\in W}(\|u\|^2+\|w\|^2)=m\|u\|^2,
\]
and similarly on $V$.  The cosine equals
$\|S_{UU}^{-1/2}S_{UV}S_{VV}^{-1/2}\|$ and is at most
\[
 \frac{xy}{m}+
 \frac{xy(1+x)(1+y)}{m^2}
 =\frac{2xy(1+xy)}{(1-x)^2(1-y)^2}.
\]
\end{proof}

\begin{proof}[Proof of Theorem~\ref{q:thm-wire-gap}]
First restrict each block to its image under $T_v$ and remove its
isometric polar factor.  This identifies the network with independent
edge vectors acted on by invertible positive block filters.  An invertible
map preserves intersections of subspaces; hence the common kernel in
these block images is exactly the image of the product of fixed internal
edge vectors, with all dangling boundary factors arbitrary.  Consider
two edges meeting at $v$, and label their ports $i,j$.  Keep the actual
attachment map $\mathcal R$ on every unselected port.  With
$D_i=\Phi_i-\Psi_i$ and $D_j=\Phi_j-\Psi_j$, take the reference metric
\[
 G_0=(\Psi_i\otimes\Psi_j\otimes\mathcal R)(C_v).
\]
It is faithful by \eqref{q:injectivity}, and has product form
$R_i\otimes R_j\otimes R_{\mathrm{rest}}$.  Whitening therefore
acts separately on the selected ports and on the remaining ones.
The filters at the two outside endpoints can be absorbed into two
reference subspaces supported on disjoint sets of factors.  Their
orthogonal projectors $p,q$ commute.  The factor acting on the retained
ports preserves both subspaces.

After whitening the exact metric has the expansion
\[
 G=G_0^{-1/2}(\Phi_i\otimes\Phi_j\otimes\mathcal R)(C_v)G_0^{-1/2}
   =I+A+B+D.
\]
Here $A$ uses $D_i\otimes\Psi_j$, $B$ uses
$\Psi_i\otimes D_j$, and $D$ uses $D_i\otimes D_j$.
The first two terms have identity on the opposite selected port, so
$[A,q]=[B,p]=0$.  The single-port sandwiches imply
$\|A\|\leq x$, $\|B\|\leq y$, for $x=\delta_i$, $y=\delta_j$.
For the mixed term use the exact CP inequalities
\[
 -xy\,\Psi_i\otimes\Psi_j
 \leq_{\mathrm{CP}}D_i\otimes D_j
 \leq_{\mathrm{CP}}xy\,\Psi_i\otimes\Psi_j.
\]
For example, the upper inequality is the average of the two CP maps
\[
 (x\Psi_i+D_i)\otimes(y\Psi_j-D_j),\qquad
 (x\Psi_i-D_i)\otimes(y\Psi_j+D_j).
\]
The other pair of signs gives the lower inequality.  Tensoring with
$\mathcal R$, applying the result to $C_v$, and whitening gives
$\|D\|\leq xy$.  Tensoring the lower single-port bounds also gives
$G\geq(1-x)(1-y)I$.  Lemma~\ref{q:lem-schur} now bounds the
Friedrichs cosine by $C_{ef}$.

Put $h_e=I-G_e$.  The two-projection decomposition implies
\begin{equation}
 \operatorname{Re}\langle h_e\xi,h_f\xi\rangle
 \geq-C_{ef}\|h_e\xi\|\|h_f\xi\|.
 \label{q:projection-cross}
\end{equation}
For completeness, the components of $h_e\xi$ and $h_f\xi$ in their
common range coincide.  Removing this component leaves vectors whose
overlap is controlled by the nontrivial principal angles.  These angles
are the same for a pair of projections and for their complements.
The removed common component contributes a nonnegative squared norm.
This proves \eqref{q:projection-cross}; commuting disjoint constraints
have zero cosine and a nonnegative cross term.

With $z_e=\|h_e\xi\|$, summing yields
\[
 \|H_{\mathrm{block}}\xi\|^2
 \geq z^{\mathsf T}(I-C)z
 \geq[1-\rho(C)]\langle\xi,H_{\mathrm{block}}\xi\rangle.
\]
The spectral theorem gives the desired gap on the product of block
images.  In the full physical space, every one-block image projector
commutes with each $G_e$.  A sector outside any block image has an
incident term equal to the identity, and hence energy at least one.
This proves the stated bound on the full Hilbert space.
\end{proof}

The matrix step is a form of the established subspace-angle method
\cite{q:Kassabov2011}.  Its representation-theoretic interpretation is
useful: averaging over a finite subgroup projects onto its fixed vectors,
and a matrix of pairwise angles controls the gap of the sum of the
corresponding averaging Laplacians.  Here CP transport supplies concrete
entries of that matrix for tensor-network subspaces.

\subsection{Microscopic normalization and decorated AKLT models}

\begin{lemma}[Finite-patch comparison]
\label{q:lem-seed}
Let $H_{\mathrm{micro}}=\sum_a h_a$ be a frustration-free microscopic
Hamiltonian with $h_a\geq0$.  For every coarse edge $e$, let $P_e$ be
its two-block patch and let $H_{P_e}$ be the sum of microscopic terms
contained in it.  Assume
\[
 \ker H_{P_e}=\mathcal R_e,\qquad
 H_{P_e}\geq g_* (I-G_e),\qquad g_*>0.
\]
Suppose every microscopic term lies in at least one and at most $m$
patches.  Then the two Hamiltonians have the same full kernel and
\begin{equation}
 \operatorname{gap}H_{\mathrm{micro}}
 \geq \frac{g_*}{m}\operatorname{gap}H_{\mathrm{block}}.
 \label{q:seed-comparison}
\end{equation}
\end{lemma}

\begin{proof}
Positivity and patch coverage give
$\ker H_{\mathrm{micro}}=\bigcap_e\ker H_{P_e}
=\bigcap_e\mathcal R_e=\ker H_{\mathrm{block}}$.
The operator inequality
$mH_{\mathrm{micro}}\geq\sum_eH_{P_e}\geq g_*H_{\mathrm{block}}$
then proves the claim above their common kernel.
\end{proof}

In the standard AKLT model \cite{q:AKLT1988}, a vertex of full graph degree $D_v$ carries
spin $D_v/2$, and every interaction is the norm-one projector onto the
maximal total spin of its two endpoints.  The spin-one wire channel is
\[
 \Phi(X)=\frac13\sum_{a=1}^3\sigma_aX\sigma_a
 =\frac23\operatorname{tr}(X)I-\frac13X.
\]
Let $\Psi(X)=\operatorname{tr}(X)I/2$.  Then
\[
 \Phi^\ell=\Psi+\beta_\ell(\mathrm{id}-\Psi),\qquad
 \beta_\ell=(-1/3)^\ell.
\]
Relative to the faithful Choi metric of $\Psi$, the Choi eigenvalues
are $1+3\beta_\ell$ once and $1-\beta_\ell$ three times.  Centering
them gives
\begin{equation}
 \delta_\ell=\frac{2}{3^\ell+(-1)^\ell},\qquad\ell\geq2.
 \label{q:aklt-radius}
\end{equation}

\begin{corollary}[Logarithmic sufficient decoration]
\label{q:cor-aklt}
Let $\mathcal G$ be a simple graph of maximum degree $\Delta\geq2$.
Replace every edge by a chain with $n=2\ell$ additional spin-one
vertices, and give each original degree-$d$ vertex spin $d/2$.
For connected finite regions made of whole half-wire blocks, with at
least two blocks, suppose
\[
 2(\Delta-1)\mathcal B(\delta_\ell,\delta_\ell)<1.
\]
Then the standard microscopic AKLT Hamiltonians have the uniform bound
\begin{equation}
 \operatorname{gap}H_{\mathrm{AKLT}}
 \geq\frac{g_*(\Delta,\ell)}\Delta
       \left[1-2(\Delta-1)
                    \mathcal B(\delta_\ell,\delta_\ell)\right]>0,
 \label{q:aklt-gap}
\end{equation}
where $g_*(\Delta,\ell)>0$ is the minimum positive gap of the finitely
many two-block patch Hamiltonians with these parameters.  In particular,
\begin{equation}
 \ell=\left\lceil\log_9(144\Delta)\right\rceil,
 \qquad n=2\ell
 \quad\Longrightarrow\quad
 \operatorname{gap}H_{\mathrm{AKLT}}
 \geq\frac{g_*(\Delta,\ell)}{2\Delta}>0.
 \label{q:log-decoration}
\end{equation}
Thus $n<\log_3(144\Delta)+2$ is a sufficient even decoration count.
\end{corollary}

\begin{proof}
The only microscopic input not yet verified is the regional kernel.
Realize spin $D_v/2$ as the homogeneous binary polynomials of degree
$D_v$ in $(x_v,y_v)$.  For an included edge $vw$, the kernel of the
maximal-spin projector consists precisely of polynomials divisible by
\[
 b_{vw}=x_vy_w-y_vx_w.
\]
Indeed, multiplication by this invariant bracket is injective and
$SU(2)$-equivariant from degrees $(D_v-1,D_w-1)$.  Its image has dimension
$D_vD_w$ and contains no maximal-spin summand.  This equals the
dimension of the complement of that summand in
$(D_v+1)(D_w+1)$ dimensions.  Distinct edge brackets are nonassociate
irreducibles in a polynomial unique factorization domain.  Simultaneous
divisibility is therefore divisibility by their product.  The remaining
degrees are exactly the free boundary virtual degrees, proving the
regional tensor-image identity.

On a half-wire patch, every boundary wire vertex has one omitted edge
and one free virtual leg.  In particular no boundary corner imposes an
additional symmetrization kernel.  The local degree and wire-length
bounds give only finitely many two-block patch types; each finite
positive matrix has a positive least nonzero eigenvalue.  Their minimum
$g_*(\Delta,\ell)$ is positive.  Internal-block microscopic terms occur
in at most $\Delta$ two-block patches, and a term crossing a block
interface occurs once.  Hence Lemma~\ref{q:lem-seed} applies with
$m\leq\Delta$.

Every edge meets at most $2(\Delta-1)$ others, proving
\eqref{q:aklt-gap}.  For the explicit choice
\eqref{q:log-decoration}, one has $\ell\geq3$,
$\delta_\ell\leq3^{1-\ell}\leq1/9$, and
\[
 \mathcal B(\delta_\ell,\delta_\ell)
 \leq4\delta_\ell^2.
\]
The last inequality follows by monotonicity and
$2(1+(1/9)^2)/(1-1/9)^4=3321/1024<4$.  Therefore
\[
 2(\Delta-1)\mathcal B(\delta_\ell,\delta_\ell)
 \leq\frac{72(\Delta-1)}{9^\ell}<\frac12.
\]
\end{proof}

The coefficient $g_*(\Delta,\ell)$ fixes the microscopic energy
normalization.  The result asserts a positive constant for each fixed
degree and decoration; it does not assert a degree-uniform physical
gap.  At small degree one can use the exact criterion rather than the
closed bound: $\ell=2$ gives $\delta=1/5$, $\mathcal B=13/64$, and
coarse gap at least $3/16$ for $\Delta=3$.  At $\ell=3$,
$\mathcal B=85/5184$, and the criterion holds through $\Delta=31$.
These values are sufficient bounds, not optimal-decoration claims.

For any fixed primitive finite-dimensional wire channel, normalized
powers converge to a faithful rank-one map.  Since its limiting Choi
matrix is positive definite, there are $C_0>0$ and $0<\theta<1$ such
that $\delta_\ell\leq C_0\theta^\ell$ for all sufficiently large
$\ell$; a Jordan polynomial is absorbed by increasing $\theta$.
The same theorem then gives sufficient total wire length
\[
 n=\frac{\log\Delta}{|\log\theta|}+O(1).
\]
This conclusion concerns the canonical blocked parent for arbitrary
nonzero cores.  A specified microscopic interaction is covered whenever
the finite regional-kernel and seed-gap hypotheses of
Lemma~\ref{q:lem-seed} hold.

\section{Compatible idempotents as local gap certificates}
\label{q:sec-oblique}

The preceding construction starts from independent edge constraints and
changes their metric.  A related algebraic description starts from
commuting, generally nonorthogonal idempotents.  Orthogonalizing their
ranges produces the physical parent Hamiltonian.  Exact compatibility
is useful because a small commutator of arbitrary orthogonal projections
does not distinguish nearly perpendicular subspaces from nearly parallel
ones.

\begin{theorem}[A compatible-idempotent certificate]
\label{q:thm-oblique}
Let $P_i^2=P_i$ and $[P_i,P_j]=0$ on a finite-dimensional Hilbert space.
Let $\Pi_i$ be the orthogonal projector onto $\operatorname{ran}P_i$,
$H=\sum_i(I-\Pi_i)$, and set
\[
 a_i=\|P_i-P_i^*\|,\qquad
 \kappa_i=\sqrt{1+a_i^2},\qquad
 L_i=\kappa_i+2\kappa_i^2a_i,\qquad
 \epsilon_{ij}=\|[P_i,P_j^*]\|.
\]
For $i\ne j$ define
\begin{equation}
 b_{ij}=\min\left\{
 \frac{a_i}{\kappa_i},\frac{a_j}{\kappa_j},
 2\min(\kappa_i,\kappa_j)L_iL_j\epsilon_{ij}\right\},
 \qquad b_{ii}=0.
 \label{q:oblique-b}
\end{equation}
If the symmetric matrix $B=(b_{ij})$ has $\rho(B)<1$, then
\[
 H\geq[1-\rho(B)](I-\Pi_{\mathcal K}),\qquad
 \mathcal K=\bigcap_i\operatorname{ran}P_i.
\]
Only the individual nonorthogonality and pairwise adjoint-commutator
norms enter the certificate.
\end{theorem}

\begin{proof}
Relative to its range and its orthogonal complement, an idempotent has
the form $P=\left(\begin{smallmatrix}I&A\\0&0\end{smallmatrix}\right)$.
Direct multiplication gives
\begin{equation}
 \Pi(P)=PP^*[I-(P-P^*)^2]^{-1},\qquad
 \|P\|\leq\sqrt{1+\|P-P^*\|^2}.
 \label{q:orthogonalizer}
\end{equation}
The denominator is at least $I$.  For any operator $T$, product and
inverse commutator identities consequently give
\begin{equation}
 \|[\Pi(P),T]\|
 \leq (\kappa+2\kappa^2a)
          \bigl(\|[P,T]\|+\|[P^*,T]\|\bigr).
 \label{q:ortho-comm}
\end{equation}
Indeed the numerator costs $\kappa$ times the bracketed sum, and the
commutator of $I-(P-P^*)^2$ costs $2a$ times that sum.  Applying
\eqref{q:ortho-comm} twice and using exact compatibility gives
\[
 \|[\Pi_i,\Pi_j]\|\leq2L_iL_j\epsilon_{ij}.
\]

To control the nearly parallel possibility, quotient by
$W=\operatorname{ran}P_i\cap\operatorname{ran}P_j$.  The two idempotents
descend with norms no larger than $\kappa_i,\kappa_j$, and their product
is zero.  For vectors $u,v$ in their respective quotient ranges,
\[
 \|u\|\leq\kappa_i\|u-tv\|\qquad(t\in\mathbb C).
\]
Minimizing in $t$ bounds their cosine by
$\sqrt{1-\kappa_i^{-2}}=a_i/\kappa_i$, and symmetrically by
$a_j/\kappa_j$.  On a nontrivial principal-angle block with cosine $c_r$,
the orthogonal-projection commutator has norm
$c_r\sqrt{1-c_r^2}$.  The preceding bounds imply
$\sqrt{1-c_r^2}\geq1/\min(\kappa_i,\kappa_j)$.
Taking the maximal cosine therefore gives
\[
 c_{ij}\leq\min(\kappa_i,\kappa_j)\|[\Pi_i,\Pi_j]\|
 \leq2\min(\kappa_i,\kappa_j)L_iL_j\epsilon_{ij}.
\]
This proves $c_{ij}\leq b_{ij}$.  The argument
\eqref{q:projection-cross} now gives $H^2\geq[1-\rho(B)]H$ and the
claimed spectral consequence.
\end{proof}

\begin{corollary}[Filtered commuting-projector models]
\label{q:cor-filtered}
Let $E_i$ be commuting local orthogonal projections and let
$L=\bigotimes_x L_x$ be any product of invertible site filters.  Then
$P_i=LE_iL^{-1}$ are local compatible idempotents.  Theorem
\ref{q:thm-oblique} applies to the physical Hamiltonian
\[
 H_L=\sum_i\bigl(I-\Pi_{L\operatorname{ran}E_i}\bigr),
 \qquad
 \ker H_L=L\bigcap_i\operatorname{ran}E_i.
\]
In particular, suppose every reference support has at most $s$ sites,
meets at most $M$ other supports, and $L_x=e^{K_x}$ with
$\|K_x\|\leq\varepsilon$.  If
\[
 2M(e^{2s\varepsilon}-1)<1,
\]
then $\operatorname{gap}H_L\geq1-2M(e^{2s\varepsilon}-1)$.
\end{corollary}

\begin{proof}
Factors of $L$ outside a reference support cancel in $P_i$, proving
locality.  Similarity preserves idempotence, commutativity, and the
intersection of the ranges.  On a support of at most $s$ sites,
integration of the similarity derivative gives
$\|P_i-E_i\|\leq e^{2s\varepsilon}-1=:\delta$.
Thus $a_i\leq2\delta$, and \eqref{q:oblique-b} gives
$b_{ij}\leq2\delta$.  Disjoint supports give $b_{ij}=0$, so the
maximum row sum is at most $2M\delta$.
\end{proof}

The common ground space in this corollary may be degenerate, including
commuting-projector code spaces.  The argument estimates a concrete
finite family of local matrices; it does not require a bound on the
condition number of the global product $L$.

\subsection{A parent from integer modular conjugates}

Let $\rho>0$ on $\bigotimes_x\mathbb C^{d_x}$ and use fixed product
bases for vectorization.  On two copies of the system let $E_x$ be the
normalized Bell projector at site $x$, and set
\[
 P_x=(\rho\otimes I)E_x(\rho^{-1}\otimes I).
\]
They are compatible idempotents with common range
$\mathbb C|\rho\rangle\!\rangle$.  For a unitary Weyl basis
$\{U_{x,a}\}_{a=1}^{d_x^2}$,
\begin{equation}
 P_x=\frac1{d_x^2}\sum_a
       (\rho U_{x,a}\rho^{-1})\otimes\overline{U_{x,a}}.
 \label{q:integer-parent}
\end{equation}
This is an exact finite-dimensional identity, so local estimates of the
integer modular conjugates give the entries of
\eqref{q:oblique-b} directly.  For example, if
$\sigma_x=I_x/d_x\otimes\operatorname{tr}_x\rho$ and
$Y_x=\rho\sigma_x^{-1}$, then for every one-site $O_x$,
\begin{equation}
 \rho O_x\rho^{-1}=Y_xO_xY_x^{-1},
 \label{q:ratio-modular}
\end{equation}
because $\sigma_x$ commutes with $O_x$.

The normalized vector $|\rho\rangle\!\rangle$ purifies
$\rho^2/\operatorname{tr}(\rho^2)$.  For a generic $\rho$ the operators
in \eqref{q:integer-parent} can have large support; locality follows
when the supplied modular conjugates are local or admit certified
decaying approximations.  This construction is related to canonical
purified parents and their mixing criteria \cite{q:LPP2025}, and to
factorizations of exact KMS parent Hamiltonians used in quantum
algorithms \cite{q:LJL2026}.  The additional input here is the explicit
compatible-idempotent certificate \eqref{q:oblique-b}.

\section{Positive modular response calculus}
\label{q:sec-kms}

For a differentiable positive matrix path, its right logarithmic
derivative need not be Hermitian.  Keeping its Hermitian and
anti-Hermitian parts together reveals a useful cancellation.  The
functional calculus below uses positive-definite scalar kernels, a
standard source of operator inequalities
\cite{q:HK1999,q:BP2000}.  We give the kernels explicitly so that their
norms, ancilla extensions, and composition rules are available without
an implicit functional-calculus convention.

Let $\rho_\lambda>0$, $K=\log\rho$, and
\[
 Q=\dot\rho\rho^{-1}=H+A=H+iJ,
 \qquad H^*=H,\quad A^*=-A,\quad J^*=J.
\]
Throughout this section powers of $\rho$ are not renormalized to have
trace one.  Renormalization adds the scalar derivative of their
normalizing factor.

\begin{theorem}[Exact response of fractional powers]
\label{q:thm-powers}
For $0\leq r\leq1$, let
$Q_r=(\partial_\lambda\rho^r)\rho^{-r}=H_r+A_r$ be its Hermitian and
anti-Hermitian decomposition.  With $\operatorname{ad}_K(O)=[K,O]$,
\begin{align}
 H_r&=\psi_r(\operatorname{ad}_K)H,
 &\psi_r(x)&=\frac{\sinh(rx)}{\sinh x},\quad\psi_r(0)=r,
 \label{q:power-H}\\
 A_r&=\chi_r(\operatorname{ad}_K)A,
 &\chi_r(x)&=\left(\frac{\sinh(rx/2)}{\sinh(x/2)}\right)^2,
 \quad\chi_r(0)=r^2.
 \label{q:power-A}
\end{align}
Both maps are completely positive, with completely bounded
operator norms $r$ and $r^2$, respectively.  Consequently
\[
 \|H_r\|\leq r\|H\|,\qquad
 \|A_r\|\leq r^2\|A\|.
\]
The estimates hold with an arbitrary ancillary system.  Moreover,
\begin{equation}
 \lim_{r\downarrow0}\frac{Q_r}{r}
 =\dot K=f(\operatorname{ad}_K)H,
 \qquad f(x)=\frac{x}{\sinh x},\quad f(0)=1.
 \label{q:log-response}
\end{equation}
\end{theorem}

\begin{proof}
Differentiation of the exponential gives
\[
 Q_r=\int_0^r e^{sK}\dot K e^{-sK}\,ds
     =\frac{e^{r\operatorname{ad}_K}-I}
                  {\operatorname{ad}_K}\dot K.
\]
Taking Hermitian and anti-Hermitian parts yields
\[
 H_r=\frac{\sinh(r\operatorname{ad}_K)}{\operatorname{ad}_K}\dot K,
 \qquad
 A_r=\frac{\cosh(r\operatorname{ad}_K)-I}
                    {\operatorname{ad}_K}\dot K.
\]
Comparison with $r=1$ proves the formulas, using
$\cosh(rx)-1=2\sinh^2(rx/2)$.  All quotients have their continuous
value at zero.

For $0<r<1$, the positive density
\begin{equation}
 p_r(t)=\frac{\sin(\pi r)}
              {2[\cosh(\pi t)+\cos(\pi r)]}
 \label{q:pr-kernel}
\end{equation}
satisfies
\begin{equation}
 \psi_r(x)=\int_{\mathbb R}p_r(t)e^{itx}\,dt,
 \qquad \int_{\mathbb R}p_r(t)\,dt=r.
 \label{q:pr-transform}
\end{equation}
One direct verification integrates
$e^{ikz}/(\cosh z+\cos\alpha)$ around a rectangle of height $2\pi i$.
The residues at $i(\pi-\alpha)$ and $i(\pi+\alpha)$ give
\[
 \int_{\mathbb R}\frac{e^{iku}}{\cosh u+\cos\alpha}\,du
 =\frac{2\pi}{\sin\alpha}
      \frac{\sinh(\alpha k)}{\sinh(\pi k)}.
\]
Evenness and continuity include $k\leq0$.  Taking $\alpha=\pi r$,
$u=\pi t$, and $k=x/\pi$ proves \eqref{q:pr-transform}.
Thus
\[
 \psi_r(\operatorname{ad}_K)(O)
 =\int_{\mathbb R}p_r(t)e^{itK}Oe^{-itK}\,dt
\]
is a positive mixture of unitary conjugations of mass $r$.  It is
completely positive with completely bounded norm $r$, attained on $I$.
Since $\chi_r(x)=\psi_r(x/2)^2$, its positive Fourier density is
$(\widetilde p_r*\widetilde p_r)(t)$, where
$\widetilde p_r(t)=2p_r(2t)$, and its total mass is $r^2$.
The endpoints are the zero and identity maps.  Finally
$p_r/r$ converges in $L^1$ to
\[
 p_0(t)=\frac\pi4\operatorname{sech}^2(\pi t/2),
\]
a probability density with Fourier transform $f$.  Alternatively the
limit in \eqref{q:log-response} follows directly by differentiating
$\rho^r=I+rK+O(r^2)$.
\end{proof}

The full response obeys the exact KMS compatibility relation
\begin{equation}
 Q^*=e^{-\operatorname{ad}_K}Q,
 \qquad A=\tanh(\operatorname{ad}_K/2)H.
 \label{q:kms-compatibility}
\end{equation}
Thus its anti-Hermitian part is constrained by its Hermitian part and
the modular generator.  It is not an independent input.  The
fractional-power maps compose coherently:
\[
 \psi_s(rx)\psi_r(x)=\psi_{rs}(x),\qquad
 \chi_s(rx)\chi_r(x)=\chi_{rs}(x).
\]
In particular, square roots give
\[
 H_{1/2}=\tfrac12\operatorname{sech}(\operatorname{ad}_K/2)H,
 \qquad
 A_{1/2}=\tfrac14\operatorname{sech}^2(\operatorname{ad}_K/4)A.
\]

\begin{proposition}[A positive rotation generator]
\label{q:prop-rotation}
Define the even function
\[
 h(x)=\frac{1-x/\sinh x}{x\tanh(x/2)},\qquad h(0)=\frac13,
 \qquad C=h(\operatorname{ad}_K)J.
\]
Then $C=C^*$ and
\begin{equation}
 \dot K=H+i[C,K],\qquad \|C\|\leq\frac13\|J\|.
 \label{q:rotation}
\end{equation}
The map $h(\operatorname{ad}_K)$ is completely positive with completely
bounded norm $1/3$.  Its Fourier density has finite exponentially
weighted mass for every weight exponent $v<2\pi$.
\end{proposition}

\begin{proof}
Differentiating the Mittag--Leffler expansion of $x\coth(x/2)$ gives
\[
 h(x)=\frac1x\frac{d}{dx}[x\coth(x/2)]
     =8\sum_{n\geq1}
       \frac{(2\pi n)^2}{[x^2+(2\pi n)^2]^2}.
\]
Its nonnegative Fourier density is
\begin{align*}
 k(t)
 &=\sum_{n\geq1}\left(\frac1{\pi n}+2|t|\right)e^{-2\pi n|t|}\\
 &=-\frac1\pi\log(1-e^{-2\pi|t|})
       +\frac{2|t|}{e^{2\pi|t|}-1}.
\end{align*}
The logarithmic singularity at zero is integrable, and its total mass
is $h(0)=1/3$.  The positive-mixture argument proves the CP and norm
claims.  For $0\leq v<2\pi$ its weighted mass is
\[
 2\sum_{n\geq1}\left[
 \frac1{\pi n(2\pi n-v)}+
 \frac2{(2\pi n-v)^2}\right]<\infty.
\]
Equation \eqref{q:kms-compatibility} and the scalar identity
$xh(x)\tanh(x/2)=1-f(x)$ give
$i[C,K]=-[1-f(\operatorname{ad}_K)]H$, proving
\eqref{q:rotation}.
\end{proof}

If $\dot U=iCU$ and $U(0)=I$, then
\[
 \frac{d}{d\lambda}(U^*KU)=U^*HU.
\]
This separates a positive path into Hermitian forcing and unitary
rotation.  The operator-norm estimates are dimension independent.
Spatial locality, however, depends on the support growth of the
conjugations $e^{itK}(\cdot)e^{-itK}$.  The next two propositions
describe precisely why ordinary norm contraction alone does not settle
that spatial question.

\section{Locality distinctions for conditional ratios and logarithms}
\label{q:sec-obstructions}

\subsection{Local conditional ratios with a three-body logarithm}

An algebraic conditional ratio and a Hermitian conditional logarithm
encode different forms of locality.  Their distinction is related to
Belavkin--Staszewski quantum Markov chains, whose structure and recovery
maps are studied in \cite{q:BCCRJ2026}.  The following small example
quantifies the distinction and persists arbitrarily close to a product
state.

\begin{proposition}[Exact local ratios, nonlocal log difference]
\label{q:prop-ratio-obstruction}
There are faithful three-qubit states with every one-site conditional
ratio supported on the closed star of that site in the path $A-B-C$,
arbitrarily close to the maximally mixed state in operator norm, such that
$\log\rho-\log\mathbb E_A\rho$ has a nonzero $ABC$ Pauli coefficient.
Here $\mathbb E_A(T)=I_A/2\otimes\operatorname{tr}_A T$.
\end{proposition}

\begin{proof}
For $0<a,b,t<1$, take
\[
 \sigma_B=\frac{I+tX_B}{2},\quad
 U_{AB}=I+aZ_AZ_B,\quad V_{BC}=I+bZ_BZ_C,
 \quad
 \rho=\frac14\sigma_B^{1/2}U_{AB}V_{BC}\sigma_B^{1/2}.
\]
The middle factors commute and are positive, so $\rho$ is faithful.
Partial traces give $\rho_B=\sigma_B$ and
\[
 \rho_{AB}=\tfrac12\sigma_B^{1/2}U_{AB}\sigma_B^{1/2},\qquad
 \rho_{BC}=\tfrac12\sigma_B^{1/2}V_{BC}\sigma_B^{1/2}.
\]
Consequently
\begin{equation}
 \rho=\rho_{AB}\rho_B^{-1}\rho_{BC},\qquad
 \rho(\mathbb E_A\rho)^{-1}
       =\sigma_B^{1/2}U_{AB}\sigma_B^{-1/2}.
 \label{q:bs-recovery}
\end{equation}
The ratio on the right, its inverse, and its matrix logarithm act on
$AB$.  The analogous ratio at $C$ acts on $BC$.

Set $a=b=s$ and put
\[
 g(u)=\frac{\operatorname{atanh}u}{u}
      =\sum_{k\geq0}\frac{u^{2k}}{2k+1},\qquad
 v=\sqrt{t^2+\frac{4s^2(1-t^2)}{(1+s^2)^2}}.
\]
In the $Z_A=z,Z_C=w$ block, $8\rho$ equals
\[
 (1+s^2zw)I+(1+s^2zw)tX_B
           +s(z+w)\sqrt{1-t^2}\,Z_B.
\]
For a positive $2\times2$ matrix $cI+u\cdot\sigma$, the vector
coefficient of its logarithm is
$[\operatorname{atanh}(|u|/c)/|u|]u$.  Hence its $X_B$ coefficient is
$tg(v)$ in the two blocks $z=w$ and $tg(t)$ in the other two.  For
$P=Z_AX_BZ_C$, the signed average gives
\begin{equation}
 \frac18\operatorname{tr}(P\log\rho)
 =\frac t2[g(v)-g(t)]
 \geq\frac{2ts^2(1-t^2)}{3(1+s^2)^2}>0.
 \label{q:three-body-coefficient}
\end{equation}
The logarithm of $\mathbb E_A\rho$ acts on $BC$, so it has zero
$P$ coefficient.

The endpoint ratio satisfies
\[
 \|\rho(\mathbb E_A\rho)^{-1}-I\|
 =a\sqrt{\frac{1+t}{1-t}},
\]
and the middle ratio obeys
\[
 \|\rho(\mathbb E_B\rho)^{-1}-I\|
 \leq t+\frac{\sqrt{1-t^2}(a+b)}{1-ab}.
\]
Taking $a=b=t\downarrow0$ proves the smallness assertion while
\eqref{q:three-body-coefficient} stays strictly positive.
\end{proof}

The same construction on qutrits uses
$Z=\operatorname{diag}(1,-1,0)$,
$X=|0\rangle\langle1|+|1\rangle\langle0|$,
$\sigma_B=(I_3+tX_B)/3$, and prefactor $1/9$ in $\rho$.
It is real symmetric and invariant under spin-one time reversal in a
Cartesian basis.  Restricting the endpoint eigenvalues to $\pm1$ gives
the same nonzero normalized three-site coefficient.  Thus the
distinction is also present in the time-reversal-invariant spin-one
setting.  It concerns strict support; the generated distance-two term
is fully consistent with an exponentially decaying support expansion.

\subsection{A CP contraction need not contract an interaction norm}

For $n$ sites and weights $w(m)\geq1$, $w(1)=1$, define the intrinsic
interaction seminorm, modulo scalars, by
\begin{equation}
 \|O\|_w=
 \inf_{O=cI+\sum_{\varnothing\ne S\subseteq[n]}O_S}
 \sup_x\sum_{S\ni x}w(|S|)\|O_S\|,
 \label{q:intrinsic-norm}
\end{equation}
where $O_S$ acts on $S$ and as identity elsewhere.  The infimum makes
this a property of the operator, rather than of a particular expansion.

\begin{proposition}[No generic same-norm modular bound]
\label{q:prop-filter-obstruction}
For every $a,b>0$ and every such sequence of weights, there are
Hermitian $K_n,H_n$ with
\[
 \|K_n\|_w\leq a+b/n,\qquad\|H_n\|_w=1,
\]
but
\begin{equation}
 \|f(\operatorname{ad}_{K_n})H_n\|_w
 \geq n\frac{ab}{a^2+b^2}[1-f(2b)]\longrightarrow\infty,
 \qquad f(x)=x/\sinh x.
 \label{q:interaction-blowup}
\end{equation}
In particular, no bound
$\|f(\operatorname{ad}_K)H\|_w\leq C(\|K\|_w)\|H\|_w$
with locally bounded $C$ holds on this full class, even near $K=0$.
\end{proposition}

\begin{proof}
Set
\[
 D_n=\sum_{j=1}^n Z_j,\qquad V_n=X_1\cdots X_n,\qquad
 e_n=\frac a{w(n)},\qquad
 K_n=\frac bnD_n+e_nV_n,
 \qquad H_n=D_n.
\]
These displayed decompositions give the upper bounds.  For $H_n$, its
spectral diameter $2n$ gives the reverse bound one, since for any
decomposition in \eqref{q:intrinsic-norm},
\[
 \tfrac12\operatorname{diam}\operatorname{spec}(O)
 \leq\sum_S\|O_S\|
 \leq n\sup_x\sum_{S\ni x}w(|S|)\|O_S\|.
\]
The subspace spanned by $|0^n\rangle,|1^n\rangle$ is reducing for
$K_n,H_n$, on which they equal $bZ+e_nX$ and $nZ$.  Put
$r_n=\sqrt{b^2+e_n^2}$.  The modular filter fixes the Bloch component
parallel to $(e_n,0,b)$ and multiplies the perpendicular component by
$f(2r_n)$.  Hence
\[
 \langle0^n|f(\operatorname{ad}_{K_n})H_n|1^n\rangle
 =\frac{nbe_n}{r_n^2}[1-f(2r_n)].
\]
Every proper-support operator has zero matrix element between these
two product vectors.  Thus every decomposition in
\eqref{q:intrinsic-norm} has a full-support term of norm at least this
matrix element.  Multiplying by $w(n)$ and using
$b\leq r_n\leq\sqrt{a^2+b^2}$ and monotonic decrease of $f$ on
$(0,\infty)$ proves \eqref{q:interaction-blowup}.
\end{proof}

This example also explains the role of KMS compatibility.  A Hermitian
input $H$ determines the compatible tangent uniquely: in a $\rho$
eigenbasis with eigenvalues $\lambda_i$,
\[
 \dot\rho_{ij}=\frac{2\lambda_i\lambda_j}{\lambda_i+\lambda_j}H_{ij}.
\]
The resulting full response is
$Q=[I+\tanh(\operatorname{ad}_K/2)]H$.  In the reducing two-dimensional
block above, its additional anti-Hermitian part is
\[
 -i\frac{ne_n}{r_n}\tanh(r_n)Y.
\]
Therefore $\|Q_n\|_w\geq na\tanh(r_n)/r_n$, which also diverges.
Proposition~\ref{q:prop-filter-obstruction} concerns an arbitrary
Hermitian input, and does not rule out spatial estimates that use a
separately controlled full compatible response.  Together with
Theorem~\ref{q:thm-powers}, it identifies this full response as the
appropriate datum for positive-path reconstruction.

\subsection{Computational verification and formalization interfaces}

The proofs above are analytic finite-dimensional arguments.  Their
supplementary computations serve two complementary roles.  Exact
rational Choi comparisons and positive-vector inequalities are
checkable certificates for concrete channels and graph types.  Floating
point calculations independently test tensor conventions and matrix
identities; they are not used to establish the universal bounds.

The supplied two-port tests compare the quotient Schur formula with
orthogonal projectors computed from the square-root metric.  In 600
deterministically seeded examples, including nearly rank-one and
ill-conditioned cores, their largest discrepancy was
$2.75\times10^{-14}$.  Separately, direct contraction of physical AKLT
star tensors at degrees $2,3,4$ and half-wire lengths $1,2$ agreed with
\eqref{q:gram-channel} within $5\times10^{-16}$ relative error.
The modular checks comprise 20 Fourier-integral comparisons, 28 matrix
response comparisons, and 20 obstruction checks.  The supplement records
the seeds, dimensions, formulas, and complete outputs.

For formal verification, the main network proof separates into CP-order
tensoring, a Schur-complement lemma, the two-projection angle inequality,
and a polynomial divisibility proof for the AKLT regional kernel.  The
modular proof separates scalar Fourier identities from finite-dimensional
functional calculus.  This division exposes a small set of reusable
lemmas; the numerical checks do not discharge their universally
quantified proof obligations.
